\section{Field Operators}

Field operators are a powerful tool in quantum mechanics, particularly in the study of many-body systems. They allow us to describe the \uline{creation and annihilation of particles in a given state}. The field operator $\hat{\Psi}(\mathbf{r})$ can be expressed as a sum over a complete set of single-particle states:
\[
    \begin{aligned}
        \hat{\Psi}(\mathbf{r})         & = \sum_{\nu} \phi_\nu(\mathbf{r}) \hat{a}_\nu,           \\
        \hat{\Psi}^\dagger(\mathbf{r}) & = \sum_{\nu} \phi_\nu^*(\mathbf{r}) \hat{a}_\nu^\dagger,
    \end{aligned}
\]
where $\phi_\nu(\mathbf{r})$ are the single-particle wavefunctions on the phase space, and $\hat{a}_\nu$ and $\hat{a}_\nu^\dagger$ are the annihilation and creation operators for the single-particle state $\nu$, respectively.\footnote{Here we use generic ladder operators \(\hat{a}\), but according to the context, we might use \(\hat{c}\) for electrons or \(\hat{b}\) for bosons, along with the respective commutation relations.} The field operator $\hat{\Psi}(\mathbf{r})$ annihilates a particle at position $\mathbf{r}$, while the field operator $\hat{\Psi}^\dagger(\mathbf{r})$ creates a particle at position $\mathbf{r}$; they are defined in such a way that they satisfy the appropriate commutation or anticommutation relations, depending on whether we are dealing with bosons or fermions.

If we have an \textbf{homogeneous system}, we can choose plane waves as our single-particle states to sum only over momentum states, moving to the \textbf{momentum space} representation. In this case, the field operators can be expressed as:
\[
    \begin{aligned}
        \hat{\Psi}(\mathbf{r}) = \frac{1}{\sqrt{\Omega}} \sum_{\mathbf{k}} e^{-i \mathbf{k} \cdot \mathbf{r}} \hat{a}_{\mathbf{k}}, \\
        \hat{\Psi}^\dagger(\mathbf{r}) = \frac{1}{\sqrt{\Omega}} \sum_{\mathbf{k}} e^{i \mathbf{k} \cdot \mathbf{r}} \hat{a}_{\mathbf{k}}^\dagger,
    \end{aligned}
\]
we are summing over all quantum numbers $\nu$, which in this case are the momentum states $\mathbf{k}$, then we are creating or annihilating particles in all of those momentum states for a precise position $\mathbf{r}$. Since this operators are based on the same annihilation and creation operators, they satisfy the same commutation or anticommutation relations, depending on whether we are dealing with bosons or fermions. For bosons, the field operators satisfy the commutation relations:
\[
    \begin{dcases}
        [\hat{\Psi}(\mathbf{r}), \hat{\Psi}^\dagger(\mathbf{r}')]  = \delta(\mathbf{r} - \mathbf{r}'), \\
        [\hat{\Psi}(\mathbf{r}), \hat{\Psi}(\mathbf{r}')] = [\hat{\Psi}^\dagger(\mathbf{r}), \hat{\Psi}^\dagger(\mathbf{r}')] = 0,
    \end{dcases}
\]
while for fermions, they satisfy the anticommutation relations:
\[
    \begin{dcases}
        \{\hat{\Psi}(\mathbf{r}), \hat{\Psi}^\dagger(\mathbf{r}')\} = \delta(\mathbf{r} - \mathbf{r}'), \\
        \{\hat{\Psi}(\mathbf{r}), \hat{\Psi}(\mathbf{r}')\} = \{\hat{\Psi}^\dagger(\mathbf{r}), \hat{\Psi}^\dagger(\mathbf{r}')\} = 0,
    \end{dcases}
\]
derived from the commutation or anticommutation relations of the ladder operators. Instead of wavefunctions, we can also express the field operators in terms of the position basis.

\subsubsection{Interacting Hamiltonian in Terms of Field Operators}

Let us now express the Hamiltonian of an interacting system in terms of the field operators. Starting from the Hamiltonian in second quantization, derived in the previous section, it can be expressed as:
\[
    \hat{H} = \sum_{\nu_a,\,\nu_b} \hat{T}_{\nu_a \nu_b} \hat{a}_{\nu_a}^\dagger \hat{a}_{\nu_b} + \frac{1}{2} \sum_{\nu_a,\,\nu_b,\,\nu_c,\,\nu_d} \hat{V}_{\nu_c \nu_d \nu_a \nu_b} \hat{a}_{\nu_c}^\dagger \hat{a}_{\nu_d}^\dagger \hat{a}_{\nu_a} \hat{a}_{\nu_b},
\]
where we are summing over all quantum numbers $\nu$. Now rewriting it in terms of the field operators is a straightforward task, since we can express the one-body and two-body operators in a more explicit way as:
\[
    \begin{aligned}
        \hat{H} & = \sum_{\nu_a,\,\nu_b} \int \mathrm{d}^3 \mathbf{r} \, \left( \Psi_{\nu_a}^*(\mathbf{r}) T (\mathbf{r}, \nabla_{\mathbf{r}}) \Psi_{\nu_b}(\mathbf{r}) \right) \hat{a}_{\nu_a}^\dagger \hat{a}_{\nu_b}                                                                                                                                                \\
                & + \frac{1}{2} \sum_{\nu_a,\,\nu_b,\,\nu_c,\,\nu_d} \left( \int \mathrm{d}^3 \mathbf{r} \, \mathrm{d}^3 \mathbf{r}' \, \Psi_{\nu_c}^*(\mathbf{r}) \Psi_{\nu_d}^*(\mathbf{r}') V(\mathbf{r} - \mathbf{r}') \Psi_{\nu_a}(\mathbf{r}) \Psi_{\nu_b}(\mathbf{r}') \right) \hat{a}_{\nu_c}^\dagger \hat{a}_{\nu_d}^\dagger \hat{a}_{\nu_a} \hat{a}_{\nu_b},
    \end{aligned}
\]
from which is now straightforward to recognize the expressions of the field operators (with \(\Psi_{\nu}(\mathbf{r})\) the single-particle wavefunctions), so that we can rewrite the Hamiltonian as:
\[
    \hat{H} = \int \mathrm{d}^3 \mathbf{r} \, \hat{\Psi}^\dagger(\mathbf{r}) T (\mathbf{r}, \nabla_{\mathbf{r}}) \hat{\Psi}(\mathbf{r}) + \frac{1}{2} \int \mathrm{d}^3 \mathbf{r} \, \mathrm{d}^3 \mathbf{r}' \, \hat{\Psi}^\dagger(\mathbf{r}) \hat{\Psi}^\dagger(\mathbf{r}') V(\mathbf{r} - \mathbf{r}') \hat{\Psi}(\mathbf{r}') \hat{\Psi}(\mathbf{r}),
\]
where the first term represents the kinetic energy and the potential energy due to external fields, while the second term represents the interaction energy between the particles. This Hamiltonian is now expressed in terms of the field operators, and it can be used to describe systems with a variable number of particles, since the field operators can create or annihilate particles in a specific state in the system.

Let us now analyze the case in which the one-body operator is given by the kinetic term, and the two-body operator is given by the Coulomb potential, which is a common case in many-body physics.

\paragraph{Free kinetic term.}
The kinetic term of the Hamiltonian can be expressed as:
\begin{itemize}
    \item in \textbf{first quantization}, the kinetic term is given by:
          \[
              \hat{T}(\mathbf{r}, \nabla_{\mathbf{r}}) = -\frac{\hbar^2}{2m} \nabla^2_{\mathbf{r}} \delta_{\sigma \sigma'},
          \]
          where we are using the plane wave basis \(\ket{\mathbf{k}, \sigma}\), such that
          \[
              \bra{\mathbf{r}}\ket{\mathbf{k}, \sigma} = \frac{1}{\sqrt{\Omega}} e^{i \mathbf{k} \cdot \mathbf{r}},
          \]
          and if we compute the matrix element of the kinetic term, we get:
          \[
              \bra{\mathbf{k}_a, \sigma_a} \hat{T}(\mathbf{r}, \nabla_{\mathbf{r}}) \ket{\mathbf{k}_b, \sigma_b} = \frac{\hbar^2 k_b^2}{2m} \delta_{\mathbf{k}_a \mathbf{k}_b} \delta_{\sigma_a \sigma_b},
          \]
          since the plane waves are eigenstates of the kinetic energy operator (the laplacian is diagonal), and the eigenvalue is given by \(\frac{\hbar^2 k^2}{2m}\). The delta functions ensure that the kinetic term only contributes when the momentum and spin states are the same.
    \item in \textbf{second quantization}, the kinetic term can be expressed as:
          \[
              \hat{T} = \sum_{\nu_a,\,\nu_b} T_{\nu_a \nu_b} \hat{a}_{\nu_a}^\dagger \hat{a}_{\nu_b} = \sum_{\mathbf{k}, \mathbf{k}^{\prime}, \sigma, \sigma^{\prime}} \bra{\mathbf{k}, \sigma} \hat{T}(\mathbf{r}, \nabla_{\mathbf{r}}) \ket{\mathbf{k}^{\prime}, \sigma^{\prime}} \hat{a}_{\mathbf{k}, \sigma}^\dagger \hat{a}_{\mathbf{k}^{\prime}, \sigma^{\prime}} = \sum_{\mathbf{k}, \sigma} \frac{\hbar^2 k^2}{2m} \hat{a}_{\mathbf{k}, \sigma}^\dagger \hat{a}_{\mathbf{k}, \sigma},
          \]
          where we have used the fact that the kinetic term is diagonal in the plane wave basis, so that the only non-zero contributions come from terms where \(\mathbf{k} = \mathbf{k}^{\prime}\) and \(\sigma = \sigma^{\prime}\). Indeed, we can see that in the last term the kinetic term is expressed as a sum over the momentum states, with the eigenvalue of the kinetic energy operator as the coefficient, and the number operator for the corresponding momentum states:
          \[
              \hat{T} = \sum_{\mathbf{k}, \sigma} \frac{\hbar^2 k^2}{2m} \hat{n}_{\mathbf{k}, \sigma}.
          \]
\end{itemize}

\paragraph{Coulomb interaction term.}
Starting directly from the second quantized expression of the Coulomb term of the Hamiltonian, we can express it as:
\[
    V = \frac{1}{2} \sum_{\sigma,\sigma^{\prime}} \sum_{\mathbf{k}_1, \mathbf{k}_2, \mathbf{k}_3, \mathbf{k}_4} \bra{\mathbf{k}_3, \sigma; \mathbf{k}_4, \sigma^{\prime}} V(\mathbf{r} - \mathbf{r}') \ket{\mathbf{k}_2, \sigma^{\prime}; \mathbf{k}_1, \sigma} \hat{a}_{\mathbf{k}_3, \sigma}^\dagger \hat{a}_{\mathbf{k}_4, \sigma^{\prime}}^\dagger \hat{a}_{\mathbf{k}_2, \sigma^{\prime}} \hat{a}_{\mathbf{k}_1, \sigma},
\]
where we are summing over all momentum and spin states, and the matrix element of the Coulomb potential can be expressed using the normalized plane wave basis as:
\[
    V = \frac{1}{2} \sum_{\sigma,\sigma^{\prime}} \sum_{\mathbf{k}_1, \mathbf{k}_2, \mathbf{k}_3, \mathbf{k}_4} \left[ \int \int \mathrm{d}^3 \mathbf{r} \, \mathrm{d}^3 \mathbf{r}' \, \frac{e^{-i \mathbf{k}_3 \cdot \mathbf{r}}}{\sqrt{\Omega}} \frac{e^{-i \mathbf{k}_4 \cdot \mathbf{r}'}}{\sqrt{\Omega}} V(\mathbf{r} - \mathbf{r}') \frac{e^{i \mathbf{k}_2 \cdot \mathbf{r}'}}{\sqrt{\Omega}} \frac{e^{i \mathbf{k}_1 \cdot \mathbf{r}}}{\sqrt{\Omega}} \right] \hat{a}_{\mathbf{k}_3, \sigma}^\dagger \hat{a}_{\mathbf{k}_4, \sigma^{\prime}}^\dagger \hat{a}_{\mathbf{k}_2, \sigma^{\prime}} \hat{a}_{\mathbf{k}_1, \sigma},
\]
and now, exploiting the form of the Coulomb potential, we can express the matrix element as:
\[
    V = \frac{1}{2} \sum_{\sigma,\sigma^{\prime}} \sum_{\mathbf{k}_1, \mathbf{k}_2, \mathbf{k}_3, \mathbf{k}_4} \frac{e^2}{\Omega^2}\left[ \int \int \mathrm{d}^3 \mathbf{r} \, \mathrm{d}^3 \mathbf{r}' \, \frac{e^{i (\mathbf{k}_1 - \mathbf{k}_3) \cdot \mathbf{r}} e^{i (\mathbf{k}_2 - \mathbf{k}_4) \cdot \mathbf{r}'}}{\vert \mathbf{r} - \mathbf{r}^{\prime} \vert} \right] \hat{a}_{\mathbf{k}_3, \sigma}^\dagger \hat{a}_{\mathbf{k}_4, \sigma^{\prime}}^\dagger \hat{a}_{\mathbf{k}_2, \sigma^{\prime}} \hat{a}_{\mathbf{k}_1, \sigma}.
\]
We can then perform the change of variables \(\bs{\xi} = \mathbf{r} - \mathbf{r}'\), so that the integral can be expressed as:
\[
    V = \frac{1}{2} \sum_{\sigma,\sigma^{\prime}} \sum_{\mathbf{k}_1, \mathbf{k}_2, \mathbf{k}_3, \mathbf{k}_4} \frac{e^2}{\Omega^2} \left[ \int \mathrm{d}^3 \mathbf{r} \, e^{i (\mathbf{k}_1 - \mathbf{k}_3 + \mathbf{k}_2 - \mathbf{k}_4) \cdot \mathbf{r}} \int \mathrm{d}^3 \bs{\xi} \, e^{i (\mathbf{k}_2 - \mathbf{k}_4) \cdot (\mathbf{r} - \bs{\xi})} \frac{1}{\vert \bs{\xi} \vert} \right] \hat{a}_{\mathbf{k}_3, \sigma}^\dagger \hat{a}_{\mathbf{k}_4, \sigma^{\prime}}^\dagger \hat{a}_{\mathbf{k}_2, \sigma^{\prime}} \hat{a}_{\mathbf{k}_1, \sigma},
\]
where we have separated the integral over \(\mathbf{r}\) and the integral over \(\bs{\xi}\). Now the equation is basically telling us to rename \(\mathbf{k}_2 - \mathbf{k}_4\) as \(\mathbf{q}\):
\[
    V = \frac{1}{2} \sum_{\sigma,\sigma^{\prime}} \sum_{\mathbf{k}_1, \mathbf{k}_2, \mathbf{k}_3, \mathbf{q}} \frac{e^2}{\Omega^2} \left[ \int \mathrm{d}^3 \mathbf{r} \, e^{i (\mathbf{k}_1 - \mathbf{k}_3 + \mathbf{q}) \cdot \mathbf{r}} \int \mathrm{d}^3 \bs{\xi} \, e^{i \mathbf{q} \cdot (\mathbf{r} - \bs{\xi})} \frac{1}{\vert \bs{\xi} \vert} \right] \hat{a}_{\mathbf{k}_3, \sigma}^\dagger \hat{a}_{\mathbf{k}_2 - \mathbf{q}, \sigma^{\prime}}^\dagger \hat{a}_{\mathbf{k}_2, \sigma^{\prime}} \hat{a}_{\mathbf{k}_1, \sigma},
\]
where we have also relabeled the indices of the creation and annihilation operators according to the change of variables. In the end we found the \textit{Fourier transform of the Coulomb potential}:
\[
    \tilde{V} = \int \mathrm{d}^3 \bs{\xi} \, \frac{e^{i \mathbf{q} \cdot \bs{\xi}}}{\vert \bs{\xi} \vert} = \frac{4\pi}{\vert \mathbf{q} \vert^2}.
\]
As we can see, the Coulomb potential is a very \textbf{long range potential}, since its transform as \(\mathbf{q} \to 0\) diverges:
\[
    \lim_{\mathbf{q} \to 0} \tilde{V} = \frac{4 \pi}{\vert \mathbf{q} \vert^2},
\]
since it is not screened by the presence of other particles in the system, and it can have a significant effect on its properties, even at large distances. In the end, integrating over \(\mathbf{r}\) we get
\[
    V = \frac{1}{2} \sum_{\sigma,\sigma^{\prime}} \sum_{\mathbf{k}_1, \mathbf{k}_3, \mathbf{q}} \frac{e^2}{\Omega^2} \Omega \, \delta_{\mathbf{k}_3, \mathbf{k}_1 + \mathbf{q}} \, \hat{a}_{\mathbf{k}_3, \sigma}^\dagger \hat{a}_{\mathbf{k}_2 - \mathbf{q}, \sigma^{\prime}}^\dagger \hat{a}_{\mathbf{k}_2, \sigma^{\prime}} \hat{a}_{\mathbf{k}_1, \sigma},
\]
where the delta function enforces the \textbf{conservation of momentum} during the interaction, so that the total momentum before and after the interaction is the same. Finally, we can sum over the delta function, so that we can express the Coulomb term as:
\[
    V = \frac{1}{2 \Omega} \sum_{\sigma,\sigma^{\prime}} \sum_{\mathbf{k}_1, \mathbf{k}_2, \mathbf{q}} \frac{4 \pi e^2}{\vert \mathbf{q} \vert^2} \hat{a}_{\mathbf{k}_1 + \mathbf{q}, \sigma}^\dagger \hat{a}_{\mathbf{k}_2 - \mathbf{q}, \sigma^{\prime}}^\dagger \hat{a}_{\mathbf{k}_2, \sigma^{\prime}} \hat{a}_{\mathbf{k}_1, \sigma}.
\]
Thus during the interaction, two particles with momenta \(\mathbf{k}_1\) and \(\mathbf{k}_2\) exchange a momentum \(\mathbf{q}\) and end up with momenta \(\mathbf{k}_1 + \mathbf{q}\) and \(\mathbf{k}_2 - \mathbf{q}\), so that the total momentum is conserved during the interaction.\footnote{The momentum conservation is enforced at the end by the delta finction \(\delta_{\mathbf{k}_1, \mathbf{k}_3 + \mathbf{q}}\).} This is a consequence of the translational invariance of the system, which implies that the total momentum is a good quantum number:
\[
    \mathbf{k}_3 + \mathbf{k}_4 = (\mathbf{k}_1 + \mathbf{q}) + (\mathbf{k}_2 - \mathbf{q}) = \mathbf{k}_1 + \mathbf{k}_2.
\]
For the spin part, we know that the \textit{Coulomb interaction does not depend on the spin of the particles},\footnote{In some kind of interactions, we have other terms including all the Pauli matrices, so spin can switch.} so that the spin is also conserved during the interaction, which is a consequence of the rotational invariance of the system.
\begin{figure}[H]
    \begin{minipage}{0.45\textwidth}
        Thus we have modeled the interaction between the particles as a two-body interaction, which then scatters two particles with momenta \(\mathbf{k}_1\) and \(\mathbf{k}_2\) into two particles with momenta \(\mathbf{k}_1 + \mathbf{q}\) and \(\mathbf{k}_2 - \mathbf{q}\), where the momentum \(\mathbf{q}\) is the momentum exchanged during the interaction. This is a very general form of the interaction, which can be applied to any system with translational invariance, and it is the starting point for many-body perturbation theory and other techniques used to study interacting systems.
    \end{minipage}
    \hfill
    \begin{minipage}{0.5\textwidth}
        \centering
        \begin{tikzpicture}[
                particle/.style={
                        draw=blue!50,
                        thick,
                        postaction={decorate},
                        decoration={markings, mark=at position 0.55 with {\arrow[scale=1.2]{stealth}}}
                    },
                interaction/.style={
                        draw=orange!90!black,
                        thick,
                        decorate,
                        decoration={snake, amplitude=1.2mm, segment length=2.5mm}
                    }
            ]
            \coordinate (V1) at (0,-1);
            \coordinate (V2) at (0,1);

            \draw[particle] (-2, -2.5) -- (V1) node[midway, left=0.2cm, text=blue!50] {$\ket{ \mathbf{k}_2, \sigma_2 }$};
            \draw[particle] (V1) -- (2, -2.5) node[midway, right=0.2cm, text=blue!50] {$\ket{ \mathbf{k}_2 - \mathbf{q}, \sigma_2 }$};

            \draw[particle] (-2, 2.5) -- (V2) node[midway, left=0.2cm, text=blue!50] {$\ket{ \mathbf{k}_1, \sigma_1 }$};
            \draw[particle] (V2) -- (2, 2.5) node[midway, right=0.2cm, text=blue!50] {$\ket{ \mathbf{k}_1 + \mathbf{q}, \sigma_1 }$};

            \draw[interaction] (V1) -- (V2) node[midway, right=0.3cm, text=orange!90!black] {$V_q = \tfrac{4\pi e^2}{q^2}$};

            \filldraw[black] (V1) circle (2pt);
            \filldraw[black] (V2) circle (2pt);
        \end{tikzpicture}
    \end{minipage}
\end{figure}

This Hamiltonian cannot be solved exactly (for more than two particles), so we need to use some approximation methods to study the properties of the system. We are so going to develop the many body perturbation theory in the next chapters, which is a powerful tool to study interacting systems, and it is based on the idea of treating the interaction as a perturbation to the non-interacting system.